\documentclass[11pt]{article}
\usepackage{amsmath,amssymb,amsthm,geometry,hyperref}
\geometry{margin=1.1in}
\newtheorem{theorem}{Theorem}[section]
\newtheorem{corollary}[theorem]{Corollary}
\newtheorem{definition}[theorem]{Definition}
\newcommand{\R}{\mathbb{R}}
\newcommand{\C}{\mathbb{C}}
\newcommand{\Z}{\mathbb{Z}}
\title{Quantum sequential growth on the full causal tree:\\
selection, feasibility geometry, and the composite route to four
dimensions}
\author{T.~DiFiore}
\date{August 2026 --- draft}
\begin{document}
\maketitle

\begin{abstract}
We study bi-normalized quantum sequential growth of causal sets: at
each branching, complex amplitudes satisfy both the coherent sum rule
$\sum_c a_c=1$ and Born completeness $\sum_c\lvert a_c\rvert^2=1$ ---
Gudder's stochastic--unitary pair, lifted from binary c-causets to the
full Rideout--Sorkin downset tree with action phases
$a\propto e^{i\varphi\,\Delta S}$.  Three groups of results.
\emph{Selection}: the full tree collapses Gudder's solution continuum
--- $\varphi=\pi/4$ is the unique feasible phase, and across 89 weight
systems the pinned phase is determined by the root gap structure.
\emph{Feasibility geometry}: a downset's phase class is determined by
two integers (width and subtop count, mod 8); feasibility of a
restricted family is exactly cone coverage plus Born reachability of
its octant spectrum (predicting every measured wall, 166/166); this
yields an expansion law and its time-reverse, a collapse-censorship
law, with one constant, while strict Rideout--Sorkin covariance and
Bell causality are machine-checked no-gos whose obstruction lies
purely in the magnitude sector --- every phase-level condition holds
exactly.  The growth Hamiltonian is $\varphi$ times the
Benincasa--Dowker action rate; the accumulated action is a causet
invariant; the width shift and curvature clock form a $\Z_8$ Weyl
pair whose squares anticommute, supporting a Jordan--Wigner fermionic
exchange algebra with an exact parity superselection rule ---
structures that exist only for the two-dimensional action, whose
weights are the unique all-even Benincasa--Dowker member.
\emph{Dimension}: the bi-normalized class is monoidal --- the tensor
of two quantum growths is a quantum growth --- and the causal product
of two grown two-dimensional factors measures four-dimensional in the
infrared, with interval statistics converging onto the exact 4D
Minkowski benchmark ($f=0.1008$ vs $0.0994$); rotational isotropy of
the light cone emerges in the many-frame limit; field quanta
(Sorkin--Johnston modes) on the composite are delocalized with
near-linear dispersion.  All structural theorems are machine-checked
in Lean~4 (80 theorems, axiom-clean); every refuted conjecture and
failed mechanism encountered en route is reported.
\end{abstract}

\section{The kernel and its arena}

\subsection{Gudder's pair, and what fixes it}
Gudder's quantum sequential growth on c-causets \cite{gudder} imposes,
per parent, $\sum_k a_k=1$ and $\sum_k\lvert a_k\rvert^2=1$; his
Theorem~2.1 exhibits the full binary solution continuum
$a=\cos\theta\,e^{i\theta}$, $b=-i\sin\theta\,e^{i\theta}$ with
$\theta$ a free coupling per node.  We adopt the pair on the full
downset tree with the action-phase ansatz
$a_D=x_{g(D)}e^{i\varphi g(D)}$, $x\ge0$, where the gap $g(D)$ is the
Benincasa--Dowker action increment of the transition.

\begin{theorem}[Continuum collapse]
On the two-dimensional weight system, the root branching (gaps
$\pm1$) forces $\cos\varphi=\tfrac{\sqrt2}{2}$, i.e.
$\varphi=\pi/4$, and a thirteen-phase scan finds no other phase
feasible at any depth.  Gudder's continuum is collapsed to its
symmetric point by the first step of the full tree.  Across the
surviving weight systems of a $140$-system scan, the pinned phase is a
function of the root gap structure alone.
\end{theorem}

\subsection{The two-integer octant formula}
Mod 8, the four-dimensional Benincasa--Dowker coefficients
$(-1,9,-16,8)$ reduce to $(-1,+1,0,0)$:

\begin{theorem}[\texttt{octant\_formula}]
If the interval weights reduce mod 8 to $(-1,+1,0,\dots)$, then for
every downset $D$,
$g(D)\equiv 1-m_0(D)+m_1(D)\pmod 8$, where $m_0$ is the number of
maximal elements of $D$ (the causal in-degree) and $m_1$ the number of
elements with exactly one element of $D$ above them.
\end{theorem}

Two integers determine every quantum phase class.  For Poisson
sprinklings the charge $q=m_0-m_1$ obeys the exact identity
$\mathbb{E}[q]=d\,\mathbb{E}[m_0]/d\ln N$ (in 2D:
$1-e^{-N}\to1$, anchored by classical record statistics), while its
mod-8 distribution is uniform: flat ensembles cover all octants.

\section{Feasibility geometry}

\begin{theorem}[Feasibility criteria;
\texttt{single\_class\_born},
\texttt{halfplane\_separation\_infeasible},
\texttt{antipodal\_pair\_reaches\_born}]
A single gap class of multiplicity $\mu$ supports the pair iff its
phase is trivial and $\mu=1$.  If the available phases lie in a closed
half-plane missing $(1,0)$, no nonnegative solution exists.  A
nonnegative coherent solution with Born mass $\le1$ plus one antipodal
phase pair upgrades to an exact solution by an explicit quadratic
root.
\end{theorem}

These criteria predict every empirically measured wall (166/166
collapse-family tests).  Consequences: (i)~an \emph{expansion law} ---
the width a new event may cap is bounded by
$w_{\max}(n)\approx n/c^*$, $c^*\approx4.5$; (ii)~its time-reverse, a
\emph{collapse censorship} law with the same constant; (iii)~geometry
selectivity --- flat (octant-covering) ensembles are unobstructed,
walls are curvature phenomena; parallel-chain regions collapse freely
at any width (solved exactly); (iv)~free growth runs at
$\sim0.8$ of its critical expansion rate.

\begin{theorem}[No-gos; \texttt{covariance\_no\_go},
\texttt{bell\_causality\_no\_go}]
Strict Rideout--Sorkin covariance (per-path amplitude equality across
isomorphic labelled causets) and classical-form (equivalently
fixed-slot quantum) Bell causality are each inconsistent with
action-phased bi-normalization --- the former by a linear
contradiction at $n=3$, the latter at the two-antichain.  Both
obstructions lie in the magnitude sector; the phases satisfy every
covariance and causality condition exactly, and the class measure is
well-defined and consistent (verified exactly to depth 5).
\end{theorem}

The no-gos explain both Gudder's restriction to uniquely-labellable
causets (necessary, not convenient) and the long resistance of the
quantum sequential growth problem: the Rideout--Sorkin template
imposes classical-probability conditions on amplitudes, which no
bi-normalized dynamics can satisfy; the conditions must live at the
class/decoherence-functional level, where this dynamics complies.

\section{Operator structure}

\begin{theorem}[Action invariance;
\texttt{growthAction\_closed\_form},
\texttt{growthAction\_iso\_invariant}]
The cumulative growth action equals
$n-\sum_{\text{related pairs}}w(\lvert I(y,x)\rvert)$ --- a pure
interval-abundance functional --- and is invariant under every order
isomorphism, for any weight system.
\end{theorem}

The Hamiltonian is $H=\varphi\,\hat G$ (the action-rate operator):
energy is curvature.  Its Born-weighted stationary state occupies the
four quadrature levels with an emergent conjugation symmetry (the
imbalance decays as $-0.20/n$).  The width shift $S$ and curvature
clock $C$ satisfy the Weyl relation $SC=\zeta CS$,
$\zeta=e^{-i\pi/4}$ (\texttt{width\_curvature\_weyl}); their squares
anticommute (\texttt{clifford\_anticommutation}); string operators
$\chi_i$ built on the growth order anticommute at distinct ranks
(\texttt{jw\_exchange}) --- Jordan--Wigner fermions whose required
ordering is supplied by growth itself.  Action parity is conserved
identically (\texttt{action\_parity}), and dynamically the $\Z_8$
charge decays exactly onto this $\Z_2$: the superselection rule
physical fermions retain.

\begin{theorem}[The 2D selection; \texttt{gap\_odd\_of\_even\_weights}]
Even weights make every gap odd.  Among Benincasa--Dowker families
only $d=2$ has all-even weights ($d=4$ is mixed-parity with a
degenerate root; $d=3$ is non-integer): the quadrature clock, Weyl
pair, Clifford sign, fermionic algebra and parity superselection
exist for the two-dimensional action alone --- a candidate mechanism
for the universal ultraviolet dimensional reduction to two.
\end{theorem}

\section{The composite route to four dimensions}

\begin{theorem}[Monoidality; \texttt{product\_bi\_normalized}]
The tensor of two bi-normalized branchings is bi-normalized; gaps
add, so the product carries the same $\pi/4$ action phase.
\end{theorem}

The causal product of two grown 2D factors is a four-dimensional
dominance (Winkler) order: at $1024$ elements its interval dimension
profile decelerates smoothly onto four
($4.38,4.24,4.14,3.98$ at $k\sim32..256$), with
$f(k{\sim}256)=0.1008$ against the exact 4D Minkowski benchmark
$0.0994$.  Dimensional additivity
$d_{\otimes}(k)=2\,d_{2D}(\sqrt k)$ holds at every checkable scale.
The polyhedral cone rounds under intersection of rotated frames
(opening angle $\to30^\circ$, the exact inscribed value; anisotropy
$\sim1/m$), the quantum ensemble is measurably rounder than any
single history ($3.0\times$), and the four-dimensional class survives
strong selection-induced coupling.  Each factor contributes an
independent fermionic species
(\texttt{species1\_anticommute}, \texttt{species\_cross\_commute},
\texttt{composite\_parity\_pair}): $k$ factors carry $\Z_2^k$.
Sorkin--Johnston field modes on the composite are delocalized over
$\sim36\%$ of the volume with a harmonic low-$k$ spectrum and
near-linear dispersion.

\section{Honest limits}

Quantum consistency does \emph{not} select the gravitational action
(deep feasibility is a broad basin over weight space; continuum
locality selects it).  The cosmological action density is
selection-band dependent.  Manifoldlikeness is not achieved: under the
theory's coherent measure two competing forces (stationary-phase
flattening, confirmed at $34\%$; extension enhancement, $2.8\times$)
admit a balance point at the sprinkling ordering fraction
($\Delta r=0.005$), and texture-targeted weights halve the interval-
texture gap, but exact class-aggregated enumeration to $n=10$
($1.6\times10^6$ classes) shows the residual texture gap growing,
not shrinking --- the entropy problem survives the coherent measure
at all tested weights.  Growth defects thermalize (no free
quasiparticles; charge-neutral pairs decay at the single-defect
rate); free excitations exist only as field modes.  Exact Lorentzian
structure is unreachable by any finite product (round cones have
infinite order dimension) and must emerge, as isotropy does, in a
many-frame limit.  No measured constant of nature is yet reproduced.

\begin{thebibliography}{9}
\bibitem{rs} D.~Rideout, R.~Sorkin, \emph{A classical sequential
growth dynamics for causal sets}, Phys. Rev. D 61, 024002 (2000).
\bibitem{gudder} S.~Gudder, \emph{An isometric dynamics for a causal
set approach to discrete quantum gravity}, arXiv:1409.3770.
\bibitem{bd} D.~Benincasa, F.~Dowker, \emph{The scalar curvature of a
causal set}, Phys. Rev. Lett. 104, 181301 (2010).
\bibitem{surya} S.~Surya, \emph{The causal set approach to quantum
gravity}, Living Rev. Relativ. 22, 5 (2019).
\bibitem{zalel} S.~Zalel, S.~Surya, \emph{A criterion for covariance
in complex sequential growth models}, Class. Quantum Grav. 37, 195002
(2020).
\bibitem{carlip} S.~Carlip, \emph{Dimensional reduction in quantum
gravity}, Class. Quantum Grav. 34, 193001 (2017).
\bibitem{sj} R.~Sorkin, \emph{Scalar field theory on a causal set in
histories form}; S.~Johnston, \emph{Feynman propagator for a free
scalar field on a causal set}, Phys. Rev. Lett. 103, 180401 (2009).
\bibitem{winkler} P.~Winkler, \emph{Random orders}, Order 1, 317
(1985).
\bibitem{p1} T.~DiFiore, \emph{Quantum mechanics from a monotone
probability, lossless time, gauge, and one interference event},
companion paper.
\end{thebibliography}
\end{document}
